\section{The fixed-order algorithm}\label{sec:fixedq}

This section decides $\mathrm{DEC}(q)$ at a prescribed order $q$. The procedure
is Algorithm~\ref{alg:fixed-q}: in each pivot chart it quantifies only over the
entries of $X$, $G$ and the certificate vector, that is over
$2q\omega+2\omega=O(q^2)$ real variables, whereas the chart coordinates of the
order-by-order procedure of \cite{fer26b} number $d(m_L-d)=(n-p-q)(q-r)$ and
still grow linearly in $n$ when $q$ is fixed; the compression of
Section~\ref{sec:compressed} is what removes that dependence.

\begin{algorithm}[!t]
\caption{Deciding $\mathrm{DEC}(q)$ for a fixed order $q$}
\label{alg:fixed-q}
\begin{algorithmic}[1]
\Require $A\in\mathbb Q^{n\times n}$, $C\in\mathbb Q^{p\times n}$,
$L_0\in\mathbb Q^{r_0\times n}$, integer $q\ge r_0$.
\Ensure \textsc{feasible} or \textsc{infeasible}.
\State $N\gets$ basis of $\ker[C;CA]$;\quad $k\gets$ number of columns of $N$
\State $h^0\gets L_0N$;\quad $g^0\gets L_0AN$;\quad
$D\gets$ rational matrix with $\ker D=\operatorname{row}[N\ AN]$
\For{$\omega=0$ \textbf{to} $\min(q,k)$}
  \For{each $J\subseteq[k]$ with $|J|=\omega$}
    \State form the formula $\Phi_{\omega,J}(X,G,x)$ over the variables
           $X,G\in\mathbb R^{q\times\omega}$ and $x\in\mathbb R^{2\omega}$:
    \Statex \quad (i) $e_\omega(X^{\mathsf T}X)>0$ (the left inverse $U$ is
           not quantified);
    \Statex \quad (ii) the rank-consistency formula of
           Lemma~\ref{lem:rankcons} for $M(X,G)b=c(X,G)$ and $B_{:,J}=I_\omega$;
    \Statex \quad (iii) $X_iB=h_i^0$ and $G_iB=g_i^0$ for $i\le r_0$;
    \Statex \quad (iv) $D[(X_iB)^{\mathsf T};(G_iB)^{\mathsf T}]=0$ for
           $i>r_0$;
    \Statex \quad (v) $Z^{\mathsf T}Zx=e_0$, with $Z$ from the maximal minors
           of $\lambda X-G$ (Proposition~\ref{prop:gram}).
    \If{$\Phi_{\omega,J}$ is real-feasible}
      \State \Return \textsc{feasible}
    \EndIf
  \EndFor
\EndFor
\State \Return \textsc{infeasible}
\end{algorithmic}
\end{algorithm}

\begin{thm}\label{thm:fixedq}
For every fixed $q$, $\mathrm{DEC}$ is decidable in deterministic polynomial
time. More precisely, Algorithm~\ref{alg:fixed-q} runs in
$\mathcal L^{O(q^2)}$ bit operations.
\end{thm}

\begin{proof}
\emph{(i) Equivalence.} Let $L=[L_0;R]$ be feasible, $H=LN$, $H_A=LAN$,
$\omega=\operatorname{rank}H$, and let $J$ be a pivot set with
\[
H=XB,\qquad B_{:,J}=I_\omega .
\]
By Lemma~\ref{lem:compression}(ii) there is $F$ with $H_A=FH$; put $G=FX$.
Then
\[
X_iB=H_i=h_i^0,\qquad G_iB=(H_A)_i=g_i^0\qquad(i\le r_0),
\]
and, since $\ker D=\mathcal T=\operatorname{row}[N\ AN]$,
\[
D[(X_iB)^{\mathsf T};(G_iB)^{\mathsf T}]=0\qquad(i>r_0).
\]
Moreover $\operatorname{rank}X=\omega$ gives $e_\omega(X^{\mathsf T}X)>0$,
Lemma~\ref{lem:rankcons} gives the rank-consistency condition for
$B_{:,J}=I_\omega$, and Proposition~\ref{prop:gram} turns \eqref{eq:5b} into
$Z^{\mathsf T}Zx=e_0$. Hence $\Phi_{\omega,J}$ holds.

Conversely, suppose $\Phi_{\omega,J}$ holds. Lemma~\ref{lem:rankcons} yields
$B$ with $B_{:,J}=I_\omega$; put
\[
H=XB,\qquad H_A=GB .
\]
The fixed-row equations give $H_i=h_i^0$ and $(H_A)_i=g_i^0$ for $i\le r_0$,
and the free-row equations place each remaining pair in $\mathcal T$, so there
are $r_i$ with
\[
(H_i,(H_A)_i)=r_i[N\ AN]\qquad(i>r_0).
\]
With $L=[L_0;R]$, $R$ the matrix with rows $r_i$, this reads $H=LN$ and
$H_A=LAN$. Since $e_\omega(X^{\mathsf T}X)>0$, $X$ has a left inverse $U$;
hence
\[
\begin{aligned}
G&=(GU)X,\\
H_A&=GB=(GU)XB=FH,\qquad F=GU,
\end{aligned}
\]
so $\operatorname{row}H_A\subseteq\operatorname{row}H$, which is
\eqref{eq:5a} by Lemma~\ref{lem:compression}(i), and $Z^{\mathsf T}Zx=e_0$ is
\eqref{eq:5b} by Proposition~\ref{prop:gram}. Therefore the disjunction
$\bigvee_{\omega,J}\Phi_{\omega,J}$ holds if and only if a feasible $L$
exists, which is the test performed by Algorithm~\ref{alg:fixed-q}.

\emph{(ii) Count.} With $\omega\le k$ and $N=2q\omega+2\omega$,
\[
\begin{aligned}
N &= 2q\omega+2\omega\le 2q^2+2q,\\
s &= O(k)\cdot O(k)+2r_0+O(1)=O(k^2),\\
D &= \max\{O(k),\,2\omega+1\}=O(k)\le O(n),\\
B &= \operatorname{poly}(n,\mathcal L),
\end{aligned}
\]
where $t\le\min(\text{rows},\text{cols})$ of $M$ is the rank index, the first
factor $O(k)$ bounds the branches and the second the polynomials of degree
$O(k)$ carried by each branch, the additive terms are the $2r_0$ fixed-row
equations, the free-row equations and the Gram and spectral certificates of
degree $2\omega+1$, and the last line is the coefficient-size lemma of
Section~\ref{sec:fpt}. Renegar's Part II bound
$B(\log B)(\log\log B)(sD)^{O(N)}$ therefore gives, per chart,
\[
B(\log B)(\log\log B)(sD)^{O(N)}
=\operatorname{poly}(n)^{O(q^2)}
=\mathcal L^{O(q^2)},
\]
because $sD=\operatorname{poly}(n)$ and $N\le 2q^2+2q$ with $q$ fixed. The
charts number
\[
N_{\mathrm{ch}}=(q+1)k^{q}\le 2^{q}k^{q}\le\mathcal L^{2q},
\]
since $k\le n\le\mathcal L$ and $q+1\le 2^{q}\le\mathcal L^{q}$ for $q\ge1$
and $\mathcal L\ge2$. Hence
\[
N_{\mathrm{ch}}\cdot\mathcal L^{O(q^2)}
\le\mathcal L^{2q}\cdot\mathcal L^{O(q^2)}
=\mathcal L^{O(q^2)} .
\]
For fixed $q$ the exponent $O(q^2)$ is a constant, so the construction size is
$k^{O(q^2)}\operatorname{poly}(\mathcal L)=\mathcal L^{O(q^2)}$ and Renegar's
bound gives the same order. The rank index $t$ is a discrete branch, not a
degree: it is counted in $s$, not folded into $D$. The degenerate cases $k=0$
and $q=r_0$ are treated in Section~\ref{sec:deg}.
\end{proof}

Table~\ref{tab:4} collects the quantities entering the bound.

\begin{table*}[!t]
\centering
\caption{Quantities entering the fixed-order bound.}\label{tab:4}
\small
\begin{tabular}{p{0.450\textwidth}p{0.450\textwidth}}
\toprule
Quantity & Bound \\
\midrule
real variables $N$ & $2q\omega+2\omega\le2q^2+2q$ \\
pivot charts & $(q+1)k^q$ \\
rank-consistency branches $t$ & $O(k)$, one per possible rank \\
number of sign conditions $s$ & $O(k^2)$ \\
maximum degree $D$ & $O(k)$ \\
coefficient bit length $B$ & $\operatorname{poly}(n,\mathcal L)$ \\
construction size & $k^{O(q^2)}\operatorname{poly}(\mathcal L)$ \\
\bottomrule
\end{tabular}
\end{table*}

Three representation details are worth recording.

\begin{enumerate}
\item \emph{Strict inequalities.} A condition $f>0$ is written as
\[
f>0\iff\exists u:\ fu^2=1,
\]
which is equivalent over the reals. There are $O(1)$ strict inequalities per
branch, so this adds $O(1)$ variables and changes no order.
\item \emph{Disjunctions.} The rank-consistency disjunction is processed branch
by branch, with one real-feasibility query per $t$ and the running times
added:
\[
t\cdot\mathcal L^{O(q^2)}=O(k)\,\mathcal L^{O(q^2)}=\mathcal L^{O(q^2)} .
\]
This is why $t$ appears in $s$ and not in $D$.
\item \emph{Elementary symmetric polynomials.} $e_j(G)$ is represented by the
coefficients of $\det(zI+G)$, computed by fraction-free elimination on the
$k\times k$ polynomial matrix $zI+G$ in $O(k^3)$ polynomial operations of
degree $O(k)$. The expansion is never written out term by term.
\end{enumerate}

\subsection{Degenerate cases}\label{sec:deg}

\begin{enumerate}
\item \emph{$k=0$}, that is, $\ker[C;CA]=0$. Then $N$ has no columns, so
\[
H=LN=0,\qquad H_A=LAN=0\qquad\text{for every }L,
\]
and both \eqref{eq:5a} and \eqref{eq:5b} hold; in particular for $L=L_0$, so
$q_{\min}=r_0$.
\item \emph{$\omega=0$ with $k>0$.} Then $H=0$, and
\[
\operatorname{rank}(\lambda H-H_A)=0\ \ \forall\lambda\in\mathbb C
\iff H_A=0,
\]
which together with $H=0$ is exactly $\eqref{eq:5a}\wedge\eqref{eq:5b}$ for
that $L$. The branch $H=H_A=0$ of the encoding is therefore neither stronger
nor weaker than the condition.
\item \emph{$q<r_0$.} Rejected by definition, since $L=[L_0;R]$ has $q$ rows
and contains the $r_0$ rows of $L_0$.
\item \emph{$r=0$.} Decided by the exact pencil test of Section~\ref{sec:zero}.
\end{enumerate}
